% ===================================================================
%  அட்டவணை எண் 2 — மனித உடல். இடைவேளை. விளையாட்டுகள்.
%
%  Traduction, faite en 2026, en tamoul d'aujourd'hui, sur l'ido de
%  texto/io. Regles de la colonne : voir 10-attavanai-01.tex.
%
%  L'ANATOMIE EST LE MEILLEUR TEST QU'UNE LANGUE PUISSE SUBIR, et le
%  tamoul le passe sans un emprunt sur une cinquantaine de mots :
%  நெற்றி le front, பொட்டு la tempe, இமை மயிர் les cils, நாசித்
%  துளை la narine, மோவாய் le menton, கன்னத் தாடி les favoris,
%  பிடரி la nuque, விலா எலும்பு la cote, முழங்கை le coude,
%  மணிக்கட்டு le poignet, நகம் l'ongle, தொடை la cuisse, முழங்கால்
%  மடிப்பு le jarret, கெண்டைக்கால் le mollet, கணுக்கால் la
%  malleole, உள்ளங்கால் la plante, குதிகால் le talon. Les cinq sens
%  aussi : பார்வை, கேள்வி, மோப்பம், சுவை, தொடுகை.
%
%  CE N'EST PAS UN HASARD QUE CE SOIT LE CORPS QUI RESISTE A
%  L'EMPRUNT. Le tableau 1 empruntait pour le mobilier de la classe
%  et nommait en tamoul ce qu'on y enseigne ; ici c'est la meme
%  frontiere vue de l'autre cote. Ce que la langue nomme depuis
%  toujours, elle le nomme encore ; ce qui est entre avec l'ecole
%  europeenne est entre avec son mot.
%
%  LES JEUX SE PARTAGENT EN TROIS, EXACTEMENT COMME AU COREEN :
%
%   * ceux que le pays joue et nomme : பம்பரம் la toupie (18),
%     கோலிக்குண்டு les billes (21), நெட்டைக் கால் les echasses (24),
%     ஒளிந்து பிடித்தல் cache-cache (25), கண்கட்டி விளையாட்டு le
%     colin-maillard (31), பட்டம் le cerf-volant (33), ஊஞ்சல் la
%     balancoire (35), மிதிவண்டி la bicyclette (37), கால்பந்து le
%     football (36), நீச்சல் la nage (40) ;
%   * ceux qui sont venus avec leur nom : டென்னிஸ் (34), குரோக்கே
%     (15), பில்போக்கே (19) — le bilboquet n'a pas plus de nom
%     tamoul que de nom coreen, et lui en fabriquer un ferait croire
%     a un jeu qu'on connait ;
%   * ceux qu'il faut decrire : le jeu de barres (38), la main
%     chaude (28), le jeu de l'ours (32), les quilles (16).
%
%  ET LA TENTATION DU JEU DE BARRES REVIENT, DANS UNE AUTRE LANGUE.
%  Le coreen avait sous la main 진놀이, deux camps qui se
%  poursuivent ; le tamoul a கபடி, qui est cela meme, et qui est en
%  outre le jeu national du pays. On ne l'ecrit pas plus que le
%  coreen n'a ecrit le sien. Ce serait remplacer le jeu du livret
%  par un autre — et il est instructif que deux langues sans le
%  moindre rapport aient offert le meme piege au meme renvoi. La
%  colonne ecrit தடுப்பு விளையாட்டு, le jeu ou l'on barre le
%  passage, qui decrit sans nommer.
%
%  LES QUILLES N'ONT PAS DE NOM TAMOUL NON PLUS, et on ne leur en
%  forge pas : la colonne ecrit குச்சிகள் (16), les batons, ce que
%  l'ido lui-meme designe, et le mouvement du jeu est dans la
%  phrase qui suit.
%
%  LE SAUT-MOUTON ET SA NOTE. Le fac-simile distingue le
%  shultro-salto (30), ou les sauteurs appuient les mains sur les
%  EPAULES du partenaire, et le dorso-salto, ou ils n'appuient que
%  sur son DOS. La colonne coreenne avait d'abord inverti les deux ;
%  celle-ci les ecrit தோள் தாண்டல் et முதுகு தாண்டல், dans cet
%  ordre-la, et la note (*) le redit.
%
%  UN SEUL RETOURNEMENT SUR LES SEPT PAIRES DE parigi.py :
%  « kegli (16) quin li abatas per bulo (17) », une relative que le
%  tamoul ne peut pas postposer. Les six autres sont des frontieres
%  de phrase et des coordinations. Meme compte qu'au coreen, au meme
%  tableau : l'enumeration ne s'inverse pas, et ce tableau n'est
%  presque qu'une enumeration — le corps, puis les jeux.
% ===================================================================

%%K t02-apar-1 apar
\VUsaut{4.0mm}
\VUtitre{15.0pt}{60}{அட்டவணை எண் 2}
\VUsaut{2.0mm}
\VUtitre{11.4pt}{0}{மனித உடல். --- இடைவேளை.}
\VUtitre{11.4pt}{0}{விளையாட்டுகள்.}

%%K t02-tit-0 sub
\VUtitre{13.2pt}{0}{மனித உடல்.}
\VUsaut{2.4mm}
\VUcentre{10.2pt}{0}{\textit{(இயற்கை அறிவியலின் எளிய பாடம் ஒன்று.)}}
\VUsaut{3.0mm}

%%K t02-tit-1 sub pos=t02-01-1
\VUpk{10.2pt}{40}{தலை.}
\VUsaut{3.0mm}

%%K t02-01-1 p
1. --- மனித உடலின் முதன்மையான பகுதிகள் : \VUgras{தலை}, \VUgras{உடல் பகுதி},
\VUgras{உறுப்புகள்}.

%%K t02-02-1 p
2. --- தலையில் இரண்டு பகுதிகள் அடங்கும் : \VUgras{மண்டையோடு}, அதற்குள்
\VUgras{மூளை} இருக்கிறது, அதை \VUgras{தலைமுடி} மூடுகிறது ; பிறகு
\VUgras{முகம்}.

%%K t02-03-1 p
3. --- நம்முடைய படத்தில் மண்டையோட்டையும் மூளையையும் பார்க்க முடியவில்லை ;
ஆனால் முகத்தின் முக்கியமான பகுதிகளை மிகத் தெளிவாகப் பார்க்கிறோம்.

%%K t02-04-1 p
4. --- இதோ முதலில் \VUgras{நெற்றி} ; அது வலிய அறிவையும், எப்போதும் இயங்கும்
சிந்தனையையும் காட்டுகிறது. \VUgras{பொட்டுகள்} மிகவும் விசாலமாக இருக்கின்றன.
\VUgras{கண்} விறுவிறுப்பாகவும் அகலமாகத் திறந்தும் இருக்கிறது. அடர்ந்த
\VUgras{புருவமும்} நீண்ட \VUgras{இமை மயிர்களும்} அதைக் காக்கின்றன.

%%K t02-05-1 p
5. --- இதோ மேலும் \VUgras{மூக்கும்} \VUgras{நாசித் துளைகளும்}. மூக்கு நமக்கு
முகர்வதற்கும் மூச்சு விடுவதற்கும் உதவுகிறது. மூக்கின் இரு பக்கங்களிலும்
\VUgras{கன்னங்கள்}, இன்னும் தள்ளி \VUgras{காதுகள்}. முகத்தின் கீழ்ப் பகுதி
\VUgras{வாயாலும்} \VUgras{மோவாயாலும்} ஆனது.

%%K t02-06-1 p
6. --- அவற்றை நாம் நன்றாகப் பார்க்க முடியவில்லை, ஏனெனில் \VUgras{மீசையும்}
\VUgras{கன்னத் தாடியும்} அவற்றை நம்மிடமிருந்து மறைக்கின்றன. ஆனால் வாயில்
இரண்டு \VUgras{உதடுகளும்}, \VUgras{மேல்தாடையும்} \VUgras{கீழ்த்தாடையும்}
அவற்றின் \VUgras{பற்களுடன்}, \VUgras{நாக்கும்} இருக்கின்றன என்பது நமக்குத்
தெரியும். வாயால் நாம் பேசுகிறோம், சாப்பிடுகிறோம். நாக்கால் உணவைச்
சுவைக்கிறோம்.

%%K t02-07-1 p
7. --- கண், காது, மூக்கு, வாய் ஆகியவை விரல்களைப் போலவே ஐந்து புலன்களின்
உறுப்புகள் : \VUgras{பார்வை}, \VUgras{கேள்வி}, \VUgras{மோப்பம்},
\VUgras{சுவை}, \VUgras{தொடுகை}.

%%K t02-08-1 p
8. --- ஆகவே தலை மனித உடலின் மிகவும் சுவாரசியமான பகுதிகளுள் ஒன்று.

%%K t02-tit-2 sub
\VUsaut{4.4mm}
\VUpk{10.2pt}{40}{உடல் பகுதி.}
\VUsaut{3.0mm}

%%K t02-09-1 p
9. --- \VUgras{கழுத்து} தலையை உடலோடு இணைக்கிறது. அதில் \VUgras{தொண்டையும்}
\VUgras{பிடரியும்} அடங்கும்.

%%K t02-10-1 p
10. --- உடலில் இன்றியமையாத உறுப்புகள் பலவும் அடங்கும். \VUgras{மார்பில்}
\VUgras{நுரையீரல்கள்} இருக்கின்றன, அவற்றால் நாம் மூச்சு விடுகிறோம் ;
\VUgras{இதயமும்} இருக்கிறது, அதுவே உயிரின் மையம். இதயம் துடிப்பதை நிறுத்தினால்
உயிருள்ள உடல் இறந்து விடுகிறது.

%%K t02-11-1 p
11. --- \VUgras{விலா எலும்புகள்} மார்பைத் தாங்குகின்றன.

%%K t02-12-1 p
12. --- உடலின் மற்ற பகுதிகள் \VUgras{பக்கம்}, \VUgras{முதுகு},
\VUgras{தோள்கள்}, \VUgras{வயிறு}. நாம் தூங்குவதற்குப் பக்கத்தில் அல்லது
முதுகில் படுக்கிறோம் ; சுமைகளைத் தோளில் சுமக்கிறோம்.

%%K t02-13-1 p
13. --- வயிற்றுக்குள் \VUgras{இரைப்பையும்} \VUgras{குடல்களும்} இருக்கின்றன.
அவை உணவைச் செரிக்க நமக்கு உதவுகின்றன.

%%K t02-tit-3 sub
\VUsaut{4.4mm}
\VUpk{10.2pt}{40}{உறுப்புகள்.}
\VUsaut{3.0mm}

%%K t02-14-1 p
14. --- நமக்கு நான்கு \VUgras{உறுப்புகள்} உண்டு : இரண்டு \VUgras{கைகள்},
இரண்டு \VUgras{கால்கள்}. கைகளால் நாம் பொருள்களை எடுக்கிறோம், பிடிக்கிறோம்,
தூக்குகிறோம், சேர்க்கிறோம் ; கால்களால் நிற்கிறோம், நடக்கிறோம், ஓடுகிறோம்.

%%K t02-15-1 p
15. --- \VUgras{தோள்} கையை உடலோடு இணைக்கிறது. மிகவும் வலிமையான ஒரு
\VUgras{தசை}, பைசெப்ஸ், கையை அசைக்கிறது ; அந்தக் கையை \VUgras{முழங்கை}
\VUgras{முன்கையோடு} இணைக்கிறது. \VUgras{மணிக்கட்டு} \VUgras{கையின்}
மூட்டாக அமைகிறது ; கை \VUgras{விரல்களிலும்} \VUgras{நகங்களிலும்}
முடிகிறது. கையின் மேல் \VUgras{நரம்பை} (மற்றும் தமனியை) நாம் பிரித்து
அறிகிறோம் ; அவற்றுள் \VUgras{இரத்தம்} ஓடுகிறது.

%%K t02-16-1 p
16. --- காலின் பகுதிகள் இவை : \VUgras{தொடை}, \VUgras{முழங்கால்},
\VUgras{முழங்கால் மடிப்பு}, \VUgras{கெண்டைக்கால்} — அதன் \VUgras{சதையை}
\VUgras{தோல்} மூடுகிறது —, \VUgras{கணுக்கால்}, \VUgras{பாதம்}. காலின்
\VUgras{எலும்புகள்} மிகவும் பருமனாகவும் மிகவும் வலிமையாகவும் இருக்கின்றன.
பாதத்தின் நுனியில் \VUgras{கால் விரல்கள்} இருக்கின்றன. மற்ற பகுதிகள்
\VUgras{உள்ளங்காலும்} \VUgras{குதிகாலும்}.

%%K t02-17-1 p
17. --- மனித உடலின் கிட்டத்தட்ட முழுமையான விளக்கம் இதுவே.

%%K t02-17-2 p
\textit{குறிப்பு.} --- \textit{« எனக்கு... (தலை முதலியவற்றில்) வலிக்கிறது »}
\textit{என்னும் சொல்லாட்சியின் துணையால், ஆசிரியர் இந்தச் சொற்கள் அனைத்தையும்}
\textit{நல்ல அல்லது கெட்ட} \VUgras{உடல் நிலை} \textit{என்னும் பார்வையில்}
\textit{மீண்டும் பயன்படுத்தலாம்.}

%%K t02-tit-4 sub
\VUsaut{5.0mm}
\VUpk{10.2pt}{40}{உடற்பயிற்சி.}
\VUsaut{2.4mm}
\VUcentre{10.2pt}{0}{\textit{(மாணவர்களுள் ஒருவனின் கூற்று.)}}
\VUsaut{3.0mm}

%%K t02-18-1 p
18. --- நெகிழ்வாகவும் சுறுசுறுப்பாகவும் நல்ல உடல் நலத்துடனும் இருப்பதற்கு
நாம் நம் கால்களையும் கைகளையும் பயிற்ற வேண்டும். அதனால்தான் நாம்
விளையாடுகிறோம், \VUgras{உடற்பயிற்சி} செய்கிறோம்.

%%K t02-19-1 p
19. --- வாரத்தில் இந்த இரண்டு பயிற்சிகளும் மேல்நிலைப் பள்ளியின் மைதானத்தில்
நடக்கின்றன (அட்டவணை எண் 1 ஐப் பார்க்க). ஆனால் வியாழனும் ஞாயிறும் நாம் பெரிய
புல்வெளிக்குப் போகிறோம் ; அங்கே ஓடுவதற்கும் விளையாடுவதற்கும் இடம் இருக்கிறது.

%%K t02-20-1 p
20. --- நம் ஆசிரியர்களும் பெற்றோரும் சில வேளைகளில் நம்முடன் அங்கே
வருகிறார்கள் ; ஆனால் பயிற்சிகளை நடத்துவது \VUgras{உடற்பயிற்சி
ஆசிரியர்}\textsuperscript{(12)}தான்.

%%K t02-21-1 p
21. --- புல்வெளியில், நுழைவாயிலின் இடப் பக்கத்தில், \VUgras{உடற்பயிற்சிக்
கருவிகள்}\textsuperscript{(1)} இருக்கின்றன. \VUgras{ஏணியால்}\textsuperscript{(2)}
நாம் ஏறுகிறோம், \VUgras{வளையங்களிலும்}\textsuperscript{(3)}
\VUgras{தொங்கு கம்பியிலும்}\textsuperscript{(4)} தலைகீழாகச் செய்கிறோம்,
\VUgras{கயிற்று ஏணியின்}\textsuperscript{(5)} அல்லது \VUgras{முடிச்சுக்
கயிற்றின்}\textsuperscript{(6)} உச்சி வரை கைகளால் நம்மை இழுத்து ஏறுகிறோம்.

%%K t02-22-1 p
22. --- அதே வேளையில் நம் தோழர்கள் \VUgras{மரக் குதிரையின்}\textsuperscript{(7)}
மேலோ \VUgras{இணைக் கம்புகளின்}\textsuperscript{(11)} மேலோ தாண்டுகிறார்கள்.
வேறு சிலர் \VUgras{குத்துச்சண்டை போடுகிறார்கள்}\textsuperscript{(8)} அல்லது
முகக் கவசம் அணிந்து \VUgras{மெல்லிய வாளால்} \VUgras{வாள்
சண்டை செய்கிறார்கள்}\textsuperscript{(9)}. கடைசியாக வேறு சிலர் \VUgras{பளு
தூக்குகிறார்கள்}\textsuperscript{(10)}.

%%K t02-tit-5 sub
\VUsaut{4.6mm}
\VUpk{10.2pt}{40}{விளையாட்டுகள்.}
\VUsaut{3.0mm}

%%K t02-23-1 p
23. --- உடற்பயிற்சி செய்பவர்களைச் சுற்றி சிறுவர்களும் சிறுமிகளும் பல
\VUgras{விளையாட்டுகளை} விளையாடுகிறார்கள். சிறுமிகள் தங்கள் \VUgras{உருட்டு
வளையத்தால்}\textsuperscript{(14)} விளையாடுகிறார்கள் ; \VUgras{கயிற்றால்}
\textsuperscript{(13)} தாண்டுகிறார்கள், அல்லது தங்கள் தம்பிகளையும்
உறவினர்களையும் \VUgras{குரோக்கே}\textsuperscript{(15)} ஆட்டத்திற்கு
அழைக்கிறார்கள். எதிராளி வளைவின் வழியாக எளிதாகக் கடந்து விடலாம் என்று
நினைக்கும் தருணத்தில், அவனுடைய \VUgras{பந்தைத்} தங்கள் \VUgras{மரச்
சுத்தியால்} தள்ளி விடுவதில் அவர்களுக்கு மகிழ்ச்சி !

%%K t02-24-1 p
24. --- சிறுவர்கள் உடலை அதிகம் அசைக்கும் விளையாட்டுகளையே விரும்புகிறார்கள்.
விருப்பத்துடன் \VUgras{குச்சிகளால்}\textsuperscript{(16)} விளையாடுகிறார்கள் ;
அவற்றைத் திறமையாக \VUgras{உருட்டுப் பந்தால்}\textsuperscript{(17)}
கவிழ்க்கிறார்கள். \VUgras{பம்பரம்}\textsuperscript{(18)},
\VUgras{பில்போக்கே}\textsuperscript{(19)} விளையாடுவதையும், \VUgras{தவளை
விளையாட்டையும்}\textsuperscript{(20)} கூட அவர்கள் இகழ்வதில்லை. ஆனால் அவர்கள்
மிகவும் விரும்பும் விளையாட்டுகள் \VUgras{கோலிக்குண்டுகளும்}
\textsuperscript{(21)} \VUgras{சுவர்ப் பந்தும்}\textsuperscript{(22)}தான் ;
ஏனெனில் \VUgras{கண்ணாடிக் கூடத்தின்}\textsuperscript{(26)} சுவர் இலேசான
பந்தைத் திருப்பி அடிக்க மிகவும் ஏற்றது.

%%K t02-25-1 p
25. --- சிறிய இயோவான்னெஸ் தன் \VUgras{நெட்டைக் கால்களால்}\textsuperscript{(24)}
நடந்து செல்கிறான் ; அவன் மிகவும் திறமையானவன், சமநிலையை நன்றாகக் காக்கத்
தெரிந்தவன். அவனுக்கு அருகில் சில தோழர்கள் அந்தக் கண்ணாடிக் கூடத்திலும்
\VUgras{வேலிக்குப்} பின்னாலும் \VUgras{ஒளிந்து பிடித்தல்}
\textsuperscript{(25)} விளையாடுகிறார்கள். வேறு சிலர் \VUgras{அடித்தவனைக்
கண்டுபிடித்தலையும்}\textsuperscript{(28)} அல்லது ஒரு \VUgras{மட்டையிலிருந்து}
இன்னொரு மட்டைக்குப் பறக்கும் \VUgras{பறக்கும் பந்தையும்}\textsuperscript{(29)}
விரும்புகிறார்கள்.

%%K t02-noto-1 noto
\VUnotes{143.2mm}{%
(*) குழந்தைகளின், சிறுவர்களின் விளையாட்டுகளுக்கு அந்தந்த நாட்டுக்கே உரிய
பெயர்கள் பெரும்பாலும் உண்டு. அவற்றை இடோ மொழியில் இயன்ற அளவு நன்றாகப்
பெயரிட முயன்றோம் ; சில வேளைகளில் அந்த விளையாட்டைக் குறிக்கும் செயலிலிருந்தே
பெயரை எடுத்தோம், சில வேளைகளில் இந்த நாட்டு அல்லது அந்த நாட்டுப் பெயர் மிகவும்
பொருந்தாததாக இல்லாவிட்டால் அதையே தேர்ந்தெடுத்தோம். எடுத்துக்காட்டு :
\VUgras{பீப்பாய் விளையாட்டு} (ஜெர்மனியிலும் பிரான்சிலும்) என்பது
\VUgras{தவளை விளையாட்டு} \textsuperscript{(20)} (ஆங்கிலத்தில்) ; அதில்
விளையாடுபவர்கள் தங்கள் வட்டமான சிறு வில்லைகளைத் துளையிட்ட பெட்டியின்
(பீப்பாயின்) மேல் வாய் பிளந்து நிற்கும் உலோகத் தவளையின் வாய் வழியாக
எறிய முயல்கிறார்கள்.
\VUgras{தோள் தாண்டல்}\textsuperscript{(30)} \VUgras{முதுகு தாண்டலிலிருந்து}
இதனால் மட்டுமே வேறுபடுகிறது : தலைக்கு மேலாகத் தாண்டுவதற்கு விளையாடுபவர்கள்
தங்கள் கைகளைத் துணைவனின் \textit{தோள்களின்} மேல் ஊன்றுகிறார்கள் ;
« \VUgras{முதுகு தாண்டலில்} » அவனுடைய \textit{முதுகின்} மேல் மட்டுமே
ஊன்றுகிறார்கள். --- அல்லது கலந்து கொள்பவர்களுள் சிலர் குனிந்து நிற்க,
மற்றவர்கள் கைகளைத் தோளின் மேல் ஊன்றி அவர்களுடைய முதுகின் மேலாகத்
தாண்டுவார்கள் ; முதலாமவர்கள் அதிர்ச்சியை வளையாமல் தாங்க வேண்டும்,
இரண்டாமவர்கள் சமநிலையை இழக்காமல் தாண்ட வேண்டும் (அட்டவணையையே பார்க்க).%
}

%%K t02-26-1 p
26. --- புல்வெளியின் நடுவில் இரண்டு மரங்கள் இருக்கின்றன. அவற்றுள் ஒன்றின்
அடியில் ஏழு அல்லது எட்டுச் சிறுவர்கள் \VUgras{தோள் தாண்டல்}
\textsuperscript{(30)} (*) ஆட்டம் ஒன்றை ஏற்பாடு செய்திருக்கிறார்கள். எல்லா
நாடுகளிலும் இந்த விளையாட்டு தெரிந்திருப்பதில்லை ; ஆனால் \VUgras{முதுகு
தாண்டல்} என்றால் என்ன என்பது எல்லாருக்கும் தெரியும் — இந்த முறை
சிறுவர்கள் அதை விளையாடவில்லை.

%%K t02-tit-6 sub
\VUsaut{5.0mm}
\VUpk{10.2pt}{40}{வெளியில் விளையாடும் விளையாட்டுகள்.}
\VUsaut{3.4mm}

%%K t02-27-1 p
27. \VUgras{கண்கட்டி விளையாட்டும்}\textsuperscript{(31)} மிகவும் மகிழ்ச்சி
தருவது ; சிறுமிகளும் சிறுவர்களும் மாறி மாறித் தங்கள் கண்களில்
\VUgras{கண்கட்டைக்} கட்டிக் கொண்டு, பிடிபடும் திறமையற்றவர்களைப்
பிடிக்கிறார்கள்.

%%K t02-28-1 p
28. --- புல்வெளியின் நடுவில், இதோ மிகவும் பிரெஞ்சுத்தனமான ஆனால் கொஞ்சம்
கடுமையான ஒரு விளையாட்டு : அது \VUgras{கரடி விளையாட்டு}\textsuperscript{(32)},
அமைதியான \VUgras{பட்டம்}\textsuperscript{(33)} விளையாட்டை விடவும், ஆங்கில
விளையாட்டான \VUgras{டென்னிஸை}\textsuperscript{(34)} விடவும் மிகவும்
ஆரவாரமானது.

%%K t02-29-1 p
29. --- இந்தக் கடைசி விளையாட்டு பெரிய மாணவர்களின் மகிழ்ச்சி. நம் ஆசிரியர்களே
விருப்பத்துடன் அதை விளையாடுகிறார்கள். அதற்கு மிகுந்த திறமையும் சுறுசுறுப்பும்
வேண்டும், அது எனக்கு மிகவும் பிடிக்கும். \VUgras{ஊஞ்சல்}\textsuperscript{(35)}
விளையாட்டுகளைச் சிறுமிகளுக்கே விட்டு விடுகிறேன்.

%%K t02-30-1 p
30. --- ஒரு வேளை ஆபத்தாகி விடக்கூடிய இன்னொரு ஆங்கில விளையாட்டு எனக்குப்
பிடிக்காது.

%%K t02-30-1b p
நான் சொல்ல வருவது \VUgras{கால்பந்தைப்}\textsuperscript{(36)} பற்றி ; அது
என் அண்ணனுக்கு மகிழ்ச்சி தருகிறது. அதே அளவு உடல் நலத்திற்கு உகந்ததும்
உண்மையில் குறைந்த கடுமையுள்ளதுமான நம் பழைய \VUgras{தடுப்பு விளையாட்டையே}
\textsuperscript{(38)} நான் விரும்புகிறேன்.

%%K t02-tit-7 sub
\VUsaut{4.6mm}
\VUpk{10.2pt}{40}{மிதிவண்டியும் பந்தாட்டமும்.}
\VUsaut{3.0mm}

%%K t02-31-1 p
31. --- \VUgras{மிதிவண்டியில்}\textsuperscript{(37)} புல்வெளியைச் சுற்றி
வருவதையும் நான் விருப்பத்துடன் செய்கிறேன்.

%%K t02-32-1 p
32. --- அடுத்த ஆண்டு \VUgras{குதிரையேற்றம்}\textsuperscript{(39)} கற்கப்
போகிறேன். அழகாக அடி வைத்து நடக்கும், விரைந்து ஓடும், தடைகளின் மேலாகப்
பாய்ந்து தாண்டும் குதிரை எவ்வளவு அழகானது !

%%K t02-33-1 p
33. --- கோடையில் நாம் \VUgras{நீச்சல்}\textsuperscript{(40)} கற்கிறோம்.
ஆற்றின் தெளிந்த குளிர்ந்த நீரில் மகிழ்ச்சியுடன் மூழ்குகிறோம், குளிக்கிறோம்.
சில வேளைகளில் கடைசியாக ஒரு காகிதப் \VUgras{பலூனை}\textsuperscript{(41)}
விடுகிறோம் ; சூடான காற்றால் அது நிரம்பியவுடன் கம்பீரமாக வானில் ஏறுகிறது.

%%K t02-34-1 p
34. --- வேறு பல விளையாட்டுகளும் உண்டு, ஆனால் அவற்றை எண்ணி முடிக்க
முடியாது ; இந்தச் சுருக்கத்திலிருந்தே, வியாழக் கிழமைகளில் நம் அழகான
புல்வெளியில் பெரிதும் சலிப்பு வருவதில்லை என்பது புரிகிறது.

%%K t02-34-2 p
கு.-கு. --- \textit{ஆசிரியர் இந்த அதிகாரத்தை எளிதாக நிறைவு செய்யலாம் ;}
\textit{மிக முக்கியமான விளையாட்டுகள் ஒவ்வொன்றையும் இன்னும் விரிவாக}
\textit{விளக்குவதன் மூலம்.}
